\section{Branch-pole return in a Heisenberg congruence tower}
\label{sec:tower}

Let \(q_m=3^m\) and consider the four-regular Cayley graphs
\[
X_m=
\operatorname{Cay}\left(
H(\Z/q_m\Z),\{x,x^{-1},y,y^{-1}\}\right).
\]
Let \(B_2\) be the two-loop bouquet labelled by \(x,y\). Each
\(X_m\to B_2\) is a regular voltage cover with deck group
\(G_m=H(\Z/q_m\Z)\). Separately, reduction modulo \(q_{m-1}\) gives the
regular level cover \(X_m\to X_{m-1}\). The Fourier/Artin decomposition over
\(G_m\) contains the adjacency block
\[
A_\rho=\rho(x)+\rho(x)^*+\rho(y)+\rho(y)^*
\]
for each irreducible unitary \(\rho\). A sector is \emph{new at level \(m\)}
if \(\rho\) does not factor through \(G_{m-1}\).

The adjacency eigenvalue \(\lambda\) of a four-regular block contributes a
Bass factor
\[
1-\lambda u+3u^2.
\]
The graph Ramanujan condition requires every nontrivial adjacency eigenvalue
to lie in \([-2\sqrt3,2\sqrt3]\). We use this graph-theoretic condition only
as a pole-circle diagnostic; it is not the classical Riemann hypothesis.

\subsection{Trivial-only subtraction already fails}

\begin{proposition}[A conductor-new abelian branch]
\label{prop:abelian-tower}
The character
\[
\chi_m(x)=e^{2\pi i/q_m},\qquad \chi_m(y)=1
\]
is new at level \(m\), and its adjacency eigenvalue is
\[
\lambda_m^{\rm ab}=2+2\cos(2\pi/q_m)\longrightarrow4.
\]
For every \(m\ge2\), \(\lambda_m^{\rm ab}>2\sqrt3\). The two roots of its
Bass factor are real and converge to \(1/3\) and \(1\).
\end{proposition}

\begin{proof}
A one-dimensional character of frequency \(a\) factors through level
\(m-1\) exactly when \(3\mid a\); frequency one is new. Summing the four
generator phases gives the displayed eigenvalue. At \(m=2\),
\(2+2\cos(2\pi/9)>2\sqrt3\), and the sequence increases to four. The roots
\[
u_{\pm,m}=
\frac{\lambda_m^{\rm ab}
\pm\sqrt{(\lambda_m^{\rm ab})^2-12}}{6}
\]
therefore converge to the roots of \(1-4u+3u^2\).
\end{proof}

Proposition~\ref{prop:abelian-tower} proves that deleting only the exact
trivial representation does not create a uniform spectral gap. One could
respond by deleting all abelian sectors. The primitive Schrödinger blocks
show that this stronger repair also fails.

\subsection{A primitive nonabelian branch}

The Schrödinger and Harper objects below are standard
\citep{dedeo2003spectra,beguin1997harper}. Our input is the
conductor-resolved Rayleigh certificate and its use as a new-sector-gap
obstruction.

\begin{lemma}[Exact-conductor Schrödinger Artin sector]
\label{lem:schrodinger-sector}
Put \(\omega_m=e^{-2\pi i/q_m}\), and on
\(\C^{\Z/q_m\Z}\) define
\[
(U_mf)(t)=f(t+1),\qquad (V_mf)(t)=\omega_m^t f(t).
\]
Then \(U_mV_m=\omega_mV_mU_m\), and
\[
\rho_m(a,b,c)=\omega_m^cV_m^bU_m^a
\]
is an irreducible unitary representation of \(G_m\). It is new at level
\(m\), and its \(B_2\)-cover adjacency block is
\[
H_m=U_m+U_m^*+V_m+V_m^*.
\]
This block remains in the level-new Artin quotient after all factors that
descend from \(G_{m-1}\) cancel.
\end{lemma}

\begin{proof}
Using \(U_m^aV_m^{b'}=\omega_m^{ab'}V_m^{b'}U_m^a\),
\[
\rho_m(a,b,c)\rho_m(a',b',c')
=\omega_m^{c+c'+ab'}V_m^{b+b'}U_m^{a+a'},
\]
which is exactly the group law
\((a,b,c)(a',b',c')=(a+a',b+b',c+c'+ab')\). Thus \(\rho_m\) is a unitary
representation with
\[
\rho_m(x)=U_m,\qquad \rho_m(y)=V_m,\qquad
\rho_m(z)=\omega_m I.
\]
An operator commuting with \(V_m\) is diagonal in its simple eigenbasis, and
commuting also with the cyclic shift \(U_m\) forces all diagonal entries to
be equal. The commutant is therefore scalar, so the representation is
irreducible. Moreover,
\[
\rho_m(z^{q_m/3})=e^{-2\pi i/3}I\ne I,
\]
whereas \(z^{q_m/3}\) lies in the kernel of \(G_m\to G_{m-1}\); hence
\(\rho_m\) does not descend.

The displayed generator images give \(A_{\rho_m}=H_m\). In the Artin
factorization for \(X_m\to B_2\), representations factoring through
\(G_{m-1}\) reproduce exactly the representation-dependent factors of
\(X_{m-1}\to B_2\) and cancel in the level quotient. Since \(\rho_m\) does
not factor, its determinant block remains.
\end{proof}

\begin{theorem}[Primitive Schrödinger branch return]
\label{thm:schrodinger}
The representation carrying \(H_m\) is new at level \(m\), and
\[
\lambda_{\max}(H_m)\longrightarrow4.
\]
Hence, for all sufficiently large \(m\), an exact-conductor nonabelian
sector violates the four-regular Ramanujan bound. The two Bass roots
associated with \(\lambda_{\max}(H_m)\) converge to \(1/3\) and \(1\).
\end{theorem}

\begin{proof}
The triangle inequality gives \(\|H_m\|\le4\). Put
\(L_m=\lfloor\sqrt{q_m}\rfloor\) and
\(n_m=2L_m+1\). For large \(m\), let \(f_m\) be the normalized indicator of
the consecutive residues \(-L_m,\ldots,L_m\). Direct calculation gives
\[
\langle f_m,(U_m+U_m^*)f_m\rangle
=2\left(1-\frac1{n_m}\right)
\]
and
\[
\langle f_m,(V_m+V_m^*)f_m\rangle
=\frac{2}{n_m}\sum_{j=-L_m}^{L_m}
\cos\frac{2\pi j}{q_m}
\ge2\cos\frac{2\pi L_m}{q_m}.
\]
The sum of the lower bounds tends to four. The Rayleigh principle and the
upper norm bound prove \(\lambda_{\max}(H_m)\to4\). At a finite level,
equality in the upper bound would force simultaneously \(U_mf=f\) and
\(V_mf=f\) for a maximizing unit vector. The Weyl relation would then give
\(f=e^{-2\pi i/q_m}f\), which is impossible. Hence
\(\lambda_{\max}(H_m)<4\) at every finite level. The Bass-root statement
follows by continuity of the quadratic roots.
\end{proof}

The threshold can be made rational. At \(q=243\), choose \(L=15\). Since
\(\cos x\ge1-x^2/2\) and \(\pi^2<10\), the Rayleigh quotient is larger than
\[
\frac{60}{31}+2-\frac{1000}{6561}
=\frac{769442}{203391},\qquad
\frac{769442}{203391}-\frac72
=\frac{115147}{406782}>0,
\qquad \frac72>2\sqrt3.
\]
Thus level five already has a certified primitive nonabelian eigenvalue
outside the Ramanujan interval. At every finite level the nontrivial
eigenvalue is strictly below four, so its roots are not the trivial roots and
cannot cancel against the topological factor. They only converge to the
branching locations.

Theorem~\ref{thm:schrodinger} rules out the uniform
Ramanujan/new-sector-gap version of the registered regular-minus-trivial
Heisenberg-tower proposal. It does not rule out a separately renormalized
infinite determinant or subtraction scheme, and it does not address
nonamenable property-\(\tau\) towers.
